\section{Conventions and preliminary results}\label{sec:preliminaries}

\subsection{Graphs, partitions and scales}

Unless explicitly stated otherwise, graphs are finite, simple and
undirected. A path is simple and has at least one edge. The
\emph{order} of a graph is its number of vertices, including any
isolated vertices that have been retained. Temporary quotients
formed by identifying vertices may have loops or parallel edges;
these are removed before any decomposition theorem is applied.

An indexed family of subgraphs partitions a graph when their
assigned edge sets are pairwise disjoint and have the required
union. Each subgraph carries only its assigned edges, rather than
all edges induced by its vertex set. Different indices may have
overlapping, or even equal, vertex sets. Order sums count those
indices separately. Whenever edges from different members are
combined, we specify the assignment and account for the edges
that remain.

For disjoint vertex sets \(A,B\), write \(E_H(A,B)\) for the set
of edges of \(H\) between them and \(e_H(A,B)=|E_H(A,B)|\).
Write \(\partial_H U\) for
the edge boundary of \(U\subseteq V(H)\). An \(r\)-vertex graph
\(H\) has \emph{cut expansion at least \(h\)} if
\begin{equation}\label{eq:cut-expansion}
 |\partial_H U|\ge h\min\{|U|,r-|U|\}
 \qquad\text{for every }U\subseteq V(H).
\end{equation}
If \(r\ge2\), singleton cuts show that its minimum degree is at
least \(h\), and consequently \(r\ge h+1\).

All logarithms denoted by \(\log\) are to base two; \(\ln\)
denotes the natural logarithm. Scales \(D\) are real, and decimal
exponents denote exact rational numbers. Inequalities comparing
integer orders with real bounds have their literal meaning.
Asymptotic notation always refers to the scale tending to
infinity. Once the explicitly fixed parameters of a lemma have
been chosen, its implicit constants and lower scale threshold
are uniform over the graphs and orders allowed by that lemma.
In the final induction the finite list of parameters is fixed
first, a common lower cutoff \(D_*\) is fixed next, and only then
is \(C\) chosen.

\subsection{A path decomposition with controlled endpoints}

We use the following form of Lov\'asz's decomposition theorem.
Disjointness in its original statement means
edge-disjointness~\cite[Theorem~1]{Lovasz1968}; see also
\cite[Theorem~21]{BM2024}.

\begin{theorem}[Lov\'asz]\label{thm:lovasz}
The edges of a finite simple graph on \(n\) vertices can be
partitioned into at most \(\floor{n/2}\) paths and cycles.
\end{theorem}

The following consequence is recorded in
\cite[Corollary~22]{BM2024}. We include the deduction because
endpoint multiplicities are used in two different constructions.

\begin{corollary}\label{cor:endpoint-paths}
Every finite simple graph on \(n\) vertices has an edge
partition into paths such that each vertex is an endpoint of at
most two paths. In particular, the partition has at most \(n\)
paths.
\end{corollary}

\begin{proof}
For \(n=0\) there is nothing to prove. Otherwise add a vertex
\(v_0\) and join it to precisely the even-degree vertices of
the original graph, including the isolated vertices. Every old
vertex now has odd degree. In a path-and-cycle partition its
number of path-endpoint occurrences must therefore be odd and
positive.

Apply \Cref{thm:lovasz} to the augmented graph, and let \(p\)
be the number of path parts. The old vertices force \(2p\ge n\),
whereas the theorem gives
\[
 \ceil{n/2}\le p
 \le \text{number of parts}
 \le \floor{(n+1)/2}=\ceil{n/2}.
\]
Thus every part is a path. Each old vertex is an endpoint
exactly once: an old vertex with at least three endpoint
occurrences, together with the other old vertices, would force
at least \(n+2\) occurrences, although \(2p\le n+1\).

Delete \(v_0\) from these paths and discard any edgeless
remnants. Every old vertex is incident to at most one deleted
edge, so it gains at most one new endpoint occurrence. The
remaining paths partition the original edge set and have
endpoint load at most two. Each positive-length path contributes
two endpoint occurrences, proving the final count.
\end{proof}

\subsection{Disjoint representatives for hypergraphs}

A hypergraph here is a finite family of nonempty subsets of a
finite ground set. A matching is a subfamily of pairwise
disjoint edges. We use the following sufficient condition,
from the hypergraph Hall theorem of Aharoni and
Haxell~\cite{AH2000}, in the bounded-rank form stated in
\cite[Theorem~6]{BM2024}.

\begin{theorem}[Aharoni--Haxell]\label{thm:aharoni-haxell}
Let \(s\) be a positive integer and let
\((\mathcal H_i)_{i\in J}\) be a finite indexed family of
hypergraphs on a common ground set. Suppose every edge is
nonempty and has size at most \(s\). If, for every nonempty
\(I\subseteq J\), the union
\(\bigcup_{i\in I}\mathcal H_i\) has a matching of size
strictly greater than \(s(|I|-1)\), then there are edges
\(E_i\in\mathcal H_i\), one for each \(i\in J\), that are
pairwise disjoint.
\end{theorem}

One may pass from the uniform-size version to this statement
by padding each edge occurrence with private fresh vertices.
Every matching before padding remains a matching after padding,
and disjoint representatives after padding have disjoint
original projections. The hypergraphs need not be distinct, and equal
edges in different indexed hypergraphs cause no difficulty.

\subsection{Elementary probabilistic estimates}

We use linearity of expectation, Markov's inequality, the union
bound, and the following binomial estimates. They are stated
explicitly to fix constants for the sampling arguments.

\begin{lemma}\label{lem:chernoff}
If \(X\) is binomial with mean \(\mu\), then
\[
 \Pr(X\le\mu/2)\le e^{-\mu/8},
 \qquad
 \Pr(X\ge2\mu)\le e^{-\mu/3}.
\]
\end{lemma}

\begin{proof}
For \(X\) a sum of independent Bernoulli variables,
\(\E e^{tX}\le\exp(\mu(e^t-1))\).
Apply Markov's inequality with \(t=-\ln2\) for the lower
tail and \(t=\ln2\) for the upper tail. The resulting exponents
are \(-\mu(1-\ln2)/2\) and \(-\mu(2\ln2-1)\),
respectively. They are at most \(-\mu/8\) and \(-\mu/3\).
\end{proof}
